\documentclass[11pt]{article}
\usepackage[margin=1in]{geometry}
\usepackage{amsmath,amssymb,amsthm}
\usepackage{booktabs}
\usepackage{array}
\usepackage{enumitem}
\usepackage[colorlinks=true,linkcolor=blue,citecolor=blue,urlcolor=blue]{hyperref}

\newcommand{\code}[1]{\texttt{#1}}
\emergencystretch=3em
\newcommand{\Unif}{\operatorname{Unif}}
\newcommand{\LN}{\operatorname{log\text{-}Normal}}
\newcommand{\ag}{\operatorname{ag}}
\newcommand{\dg}{\operatorname{dg}}
\newcommand{\sex}{g}
\newcommand{\Tentry}{T_{\mathrm{entry}}}
\newcommand{\Tmal}{T_{\mathrm{mal}}}
\newcommand{\Treg}{T_{\mathrm{reg}}}
\newcommand{\Tdist}{T_{\mathrm{dist}}}
\newcommand{\Tdiagn}{T_{\mathrm{diagn}}}
\newcommand{\Tdo}{T_{d\_\mathrm{other}}}
\newcommand{\Tdl}{T_{d\_\mathrm{lung}}}
\newcommand{\Tdeath}{T_{\mathrm{death}}}
\newcommand{\Tpred}{T_{\mathrm{pred}}}
\newcommand{\Tfinal}{T_{\mathrm{final}}}
\newcommand{\Vreg}{V_{\mathrm{reg}}}
\newcommand{\Vdist}{V_{\mathrm{dist}}}
\newcommand{\Vdiagn}{V_{\mathrm{diagn}}}
\newcommand{\ns}{\mathrm{ns}}
\newcommand{\sm}{\mathrm{s}}

\title{The MILC natural-history model:\\ mathematics, archival sampler, and the fastMILC implementation}
\author{fastMILC working notes}
\date{Draft of 15 November 2026}

\begin{document}
\maketitle

\section{Scope and sources}

This note states the MILC (MIcrosimulation Lung Cancer) natural-history model
as a set of random variables and the formulas that generate them. For each
component we give (a) the model as a probability statement, (b) the sampler in
the archival R package \code{MILC} 1.0 \cite{milc2014} (files under
\code{archival\_code/R/}), and (c) the quirks of that sampler.
Section~\ref{sec:fastmilc} summarizes how \code{fastMILC} samples each component
and where it departs from the archival sampler.

We distinguish three things throughout: the model (a distribution), the archival
sampler (an algorithm that draws from an approximation of that distribution),
and the \code{fastMILC} sampler.

\paragraph{Sources and notation.} The primary description of the model is
Chrysanthopoulou (2017) \cite{ijm2017}, cited below as IJM. We use its notation
where it has one (the TSCE parameters, $q(t)$, the Gompertz and log-normal
parameters, and the names $\Tmal$, $\Treg$, $\Tdist$, $\Tdiagn$,
$d_{\mathrm{diagn}}$, $\Tdeath$). Where IJM has no symbol, we use the
archival variable name ($\Tentry$, $\Tpred$, $\Tfinal$) or introduce one
($\zeta$, $z(t)$, $\psi$, $\rho$, $\kappa$, $\ag$, $\dg$, $F_k$, $G_k$).
IJM uses $\theta$ for the whole model-parameter vector, so we use $\psi$ for the
TSCE parameter set. Chrysanthopoulou et al.\ (2021) \cite{mdm2021} describe the
calibration of $(m,\mu_{\mathrm{reg}},\mu_{\mathrm{dist}},\mu_{\mathrm{diagn}})$.
Nanayakkara et al.\ (2024) \cite{nz2024} describe a modified variant (hazard
ratios by years since quitting; mortality not stratified by smoking) and are
cited only for parameter comparisons. Section~\ref{sec:discrepancies} lists
where IJM and the archival code disagree.

\section{The fastMILC implementation}\label{sec:fastmilc}

This section uses symbols defined in Sections~\ref{sec:notation}--\ref{sec:lc}.

\paragraph{Reference implementation.} The archival code is the reference
implementation. Where IJM differs from it, in a number or in model structure,
\code{fastMILC} follows the archival code; Section~\ref{sec:discrepancies}
records each difference. Each component is drawn for all individuals at once
over a \code{data.table} population, with one L'Ecuyer random-number stream per
component, using the samplers of \code{nhppp} (version 1.1.0).

\paragraph{Maximum age.} \code{fastMILC} simulates ages up to $M=100$ years (package
option \code{milc.max\_age}). The archival code uses 110 as the upper age of the
other-cause tail and searches for $\Tmal$ on $[0,150]$; \code{fastMILC} uses $M$
in both places. This changes the distribution of $\Tdo$ in the tail (below) and
truncates all simulated ages at $M$.

\paragraph{Population recipe.} The canonical recipe is \code{synthetic\_nlst\_recipe}:
synthetic NLST-like data with one row per person (26{,}730 current and former
smokers aged 49--79 at study entry on 2002-06-30). Each person has a birth
moment (\code{inception\_moment}; onset is simulated from birth), a study
entry moment ($\Tentry$ is the age at entry), $\zeta$, $q$, and moments of
starting and quitting smoking, which give $t_s$ and $t_q$. The onset hazard
switches at $t_s$ and $t_q$ (Section~\ref{sec:onset}), so these moments are
needed; the mortality tables do not use them. The source data have no smoking
history, so the moments are currently synthetic: $t_s$ is drawn from a
Normal(17, 4) truncated to $[10,\ t_q - 1]$ (former) or $[10,\ \Tentry - 1]$
(current), and $\Tentry - t_q \sim \Unif(0, 15)$ for former smokers, which
respects the NLST criterion of quitting within 15 years before entry. The
criterion of at least 30 pack-years is not imposed. These draws are
placeholders: the moments of starting and quitting should be obtained from the
CISNET Smoking History Generator, conditional on sex, birth cohort, and
smoking status and intensity at entry.

\paragraph{Onset (Section~\ref{sec:onset}).} $\Tmal$ is drawn from birth, not
conditioned on $\Tmal>\Tentry$, with \code{vdraw\_cumulative\_intensity}
(\code{report\_first\_K = 1}) on $[0,M]$ and the closed-form $\Lambda_{\mathrm{mal}}$
of \eqref{eq:lambdaM}. $\Lambda_{\mathrm{mal}}$ has no closed-form inverse, so the
inverse is computed by Newton's method, vectorized over all individuals. For
each individual we select the smoking segment $j$ by comparing the target $z$
with $\Lambda_{\mathrm{mal}}(t_s)$ and $\Lambda_{\mathrm{mal}}(t_q)$, form the local
target $z' = z - \Lambda_{\mathrm{mal}}(b_j) + \tilde H_{\psi_j}(b_j)$ with
\[
\tilde H_\psi(t) = H_\psi(t)-H_\psi(0)
 = K\Big[-t + B^{-1}\operatorname{log1p}\big\{B(e^{ct}-1)/c\big\}\Big],
\]
which avoids the cancellation in \eqref{eq:Htsce}, and iterate
$t \leftarrow t - \{\tilde H_{\psi_j}(t)-z'\}/h_{\psi_j}(t)$ from the segment's
right end. $\tilde H_\psi$ is increasing and convex, so these iterates decrease
monotonically to the root; rows not converged after an iteration cap (roots near
$t=0$, where $h\approx 0$) fall back to vectorized bisection. This is the same
inverse-CDF map as the archival \code{t\_mal}. No event in $[0,M]$ gives
\code{NA}, where the archival code continues the search to 150 and returns 150
if there is no root.

\paragraph{Progression (Section~\ref{sec:progression}).} The volume thresholds
are drawn exactly from the truncated law on $\mathcal A$ \eqref{eq:truncset}:
$\Vdiagn$ by inverse CDF of its truncated log-normal, $(\Vreg,\Vdist)$ by
vectorized rejection that redraws only rejected rows. As in the archival code,
$\mu_\bullet=\mathrm{sd}_\bullet$, and $\Vreg$ is drawn with the parameter vector
named \code{cdist} and $\Vdist$ with the one named \code{creg}.

\paragraph{Mortality tables (Sections~\ref{sec:oc}--\ref{sec:lc}).} Each table
$F$ ($F_k$ or $G_k$), with knots $\tau_0<\dots<\tau_J$ and $\bar F = 1-F$, is
converted to a piecewise-constant intensity on its own irregular knots,
\begin{equation}\label{eq:steplambda}
\lambda_j = \frac{\log \bar F(\tau_j) - \log \bar F(\tau_{j+1})}{\tau_{j+1}-\tau_j},\qquad j=0,\dots,J-1,
\end{equation}
and sampled with \code{vdraw\_sc\_step} (\code{report\_first\_K = 1}), once per stratum
with that stratum's knot vector. The first arrival $\tilde\tau$ satisfies
$\Pr(\tilde\tau\le\tau_j) = 1-\bar F(\tau_j)/\bar F(\tau_0)$ at every knot,
exactly. Each cause has its own independent latent time, as in the archival
code, and $\Tdeath=\min(\Tdo,\Tdl)$. Further details:
\begin{itemize}[nosep]
\item The lung atom $G_k(0)$ is merged into the first cell,
  $\lambda_0 = -\log\{1-G_k(\tau_1)\}/\tau_1$ on $[0,\tau_1)$, so $G_k$ is
  reproduced at every knot from $\tau_1$ on. The archival lookup also does not
  produce deaths at zero duration. The atom is probably death within the month of
  diagnosis (SEER survival months $=0$); we have not verified this.
\item In a terminal cell with $\bar F(\tau_J)=0$ (lung ``distant NA male'';
  other-cause ``current, >85, female, 21+''), \eqref{eq:steplambda} is infinite;
  individuals who reach that cell get a uniform time within it, as in the archival
  code.
\item Other causes: the clock is time since entry and the table is fixed by the
  age group at entry. With probability $1-F_k(\tau^{\max}_k)$ there is no event
  in the window, and $\Tdo$ is drawn from the archival tail formula with $M$ in
  place of 110, $\Unif(\Tentry,\,M-\tau^{\max}_k)$. An option gives the
  presumably intended $\Unif(\Tentry+\tau^{\max}_k,\,M)$ as a model variant. The default is the
  archival rule because $(m,\mu_{\mathrm{reg}},\mu_{\mathrm{dist}},\mu_{\mathrm{diagn}})$
  were calibrated under it. The two differ substantially: for never smokers aged
  40--44 (female) the tail applies to 98.5\%, and for entry at 42 the mean ages
  at other-cause death are 63 and 79 (68 under the archival tail with 110).
  With the archival formula, entry after $M-\tau^{\max}_k$ (84 for 16-year
  tables) gives $\Tdo<\Tentry$ in the tail, as entry after 94 does in the
  archival code; these records are excluded in Section~\ref{sec:assembly}.
\item Lung cancer: death is possible after diagnosis at any stage; the stratum
  is stage $\times$ age group at diagnosis $\times$ sex, with the literal ``NA''
  fallback strata; the residual mass $1-G_k(\tau^{\max}_k)$ is a point mass of
  lung-cancer death at $\Tdiagn+\tau^{\max}_k/12$, as in the archival code.
\end{itemize}

\paragraph{Departures from the archival sampler.} Within the tabulated windows
\code{fastMILC} draws from the tables exactly, while the archival nearest-value
lookup lags them by about one cell (Section~\ref{sec:oc}): about 0.01 in
cumulative probability early in follow-up and under 0.002 after 8 years. Between
knots the \code{fastMILC} density is piecewise exponential and the archival
density piecewise uniform; the two differ by $O(\lambda_j\Delta_j)$ within a
cell. Apart from the maximum age $M$, these are the only intended differences
in distribution from the archival code.

\section{Notation and covariates}\label{sec:notation}

All event times are ages in years on the attained-age scale (time since birth),
unless marked as a duration with~$\tau$. Quantities refer to one individual; we
omit the subscript~$i$.

\begin{center}
\small
\begin{tabular}{@{}lll@{}}
\toprule
Symbol & Meaning & Archival name \\
\midrule
$\Tentry$ & age at entry & \code{T\_entry}, \code{age\_curr} \\
$\sex \in \{\text{male},\text{female}\}$ & sex & \code{gender} \\
$\zeta \in \{\text{never},\text{former},\text{current}\}$ & smoking status at entry & \code{status} \\
$t_s,\ t_q$ & ages at starting, quitting smoking; $t_q=\infty$ if current & \code{ts}, \code{tq} \\
$q$, $q(t)$ & cigarettes per day while smoking; at age $t$ (IJM) & \code{d} \\
$\Tmal$ & age at onset of the first malignant cell & \code{T\_mal} \\
$\Treg,\ \Tdist$ & ages at onset of regional, distant stage & \code{T\_reg}, \code{T\_dist} \\
$\Tdiagn$, $d_{\mathrm{diagn}}$ & age and tumor diameter at diagnosis & \code{T\_diagn}, \code{D\_diagn} \\
$\mathrm{stage}$ & stage at diagnosis, $\in\{\text{localized},\text{regional},\text{distant}\}$ & \code{stage} \\
$\Tdo,\ \Tdl$ & latent ages at death, other causes and lung cancer & \code{T\_do}, \code{T\_dl} \\
$\Tdeath$ & age at death, $\min(\Tdo,\Tdl)$ & \code{T\_death} \\
$\Delta_{\mathrm{pred}}$; $\Tpred$ & prediction period (years); $\Tentry+\Delta_{\mathrm{pred}}$ & \code{pred\_yrs}; \code{T\_pred} \\
\bottomrule
\end{tabular}
\end{center}
IJM writes $T_{d\_\mathrm{other}}$ and $T_{d\_\mathrm{lung}}$ in its simulation
algorithm (Table A.1), and $q(t)$ for the intensity at age $t$. In the archival
model the intensity is constant while smoking, so $q(t) = q$ on $[t_s,t_q)$.

\paragraph{Stratifying functions.} The mortality tables are indexed by two
grouping functions. The age-group function is
\[
\ag(x) =
\begin{cases}
\text{``$<40$''} & x < 40,\\
\text{``}5j\text{--}(5j+4)\text{''} & 5j \le x < 5j+5,\quad j=8,\dots,16,\\
\text{``$>85$''} & x \ge 85.
\end{cases}
\]
The values of $\ag$ are the archival string labels, which must match the keys of
the mortality tables; the label ``>85'' denotes the band $x\ge 85$. In the
archival \code{age\_grp}, $x=120$ exactly returns \code{NA} (an empty
\code{which()} gives index \code{Inf}), which routes that age to the ``NA''
fallback strata of Section~\ref{sec:lc}. The intensity-group function is
$\dg(q)=$ ``1--10'' for $1\le q<11$, ``11--20'' for $11\le q<21$, ``21+'' for
$q\ge 21$ (and the label ``NA'' for $q<1$).

\section{Onset of the first malignant cell: TSCE model}\label{sec:onset}

\subsection{Model}

$\Tmal$ is the first arrival time of a nonhomogeneous Poisson process on
$[0,\infty)$ (age from birth) with intensity $\lambda_{\mathrm{mal}}(t)$, so that
\[
\Pr(\Tmal > t) = \exp\{-\Lambda_{\mathrm{mal}}(t)\}, \qquad
\Lambda_{\mathrm{mal}}(t)=\int_0^t \lambda_{\mathrm{mal}}(u)\,du .
\]

\paragraph{Constant-parameter hazard.} IJM (Eq.~1, Table~1) uses the two-stage
clonal expansion (TSCE) hazard. For a parameter set
$\psi = (\nu,\alpha,\beta,\mu)$ (normal-cell initiation rate; division, apoptosis
and malignant-conversion rates of initiated cells) and $X$ normal stem cells,
define
\[
\gamma = \alpha-\beta-\mu, \qquad
B = \tfrac12\left(-\gamma + \sqrt{\gamma^2 + 4\alpha\mu}\right), \qquad
c = \gamma + 2B, \qquad
K = \frac{\nu\mu X}{\gamma + B}.
\]
The hazard is
\begin{equation}\label{eq:htsce}
h_\psi(t) = \frac{\nu\mu X\,(e^{ct}-1)}{\gamma + B\,(e^{ct}+1)}
          = \nu\mu X\,\frac{e^{ct}-1}{\,c + B\,(e^{ct}-1)\,},
\end{equation}
where the first form is IJM Eq.~1 and the second rewrites its denominator,
$\gamma+B(e^{ct}+1) = c + B(e^{ct}-1)$. The archival code parametrizes $\psi$ by
$(\nu,\alpha,\gamma,\mu)$ and leaves $\beta=\alpha-\mu-\gamma$ implicit. Its
function \code{HT\_mal} returns
\begin{equation}\label{eq:Htsce}
H_\psi(t) = K\left[-t + \frac1B \log\!\left(\gamma + B + B e^{ct}\right)\right],
\end{equation}
and $H_\psi'(t) = K\left[-1 + c\,e^{ct}/(\gamma+B+Be^{ct})\right] = h_\psi(t)$,
because $c e^{ct} - (\gamma+B+Be^{ct}) = (\gamma+B)(e^{ct}-1)$. We also checked
this against a numerical derivative of the archival \code{HT\_mal} (agreement to six
significant digits over ages 1--85). Three properties follow from
\eqref{eq:htsce} and are used later:
\begin{enumerate}[nosep]
\item $h_\psi(0)=0$; but $H_\psi(0) = (K/B)\log c \neq 0$, so only differences
  $H_\psi(b)-H_\psi(a)$ are cumulative intensities. (The archival code uses
  differences, via \code{ht\_mal\_int}.)
\item $h_\psi$ is strictly increasing in $t$: with $y=e^{ct}-1\ge 0$,
  $y/(c+By)$ is increasing because $c>0$.
\item $h_\psi(t) \uparrow \nu\mu X / B$ as $t\to\infty$.
\end{enumerate}
$H_\psi$ has no closed-form inverse: in $y=e^{ct}$, $H_\psi$ is a linear
combination of $\log y$ and $\log(\gamma+B+By)$.

\paragraph{Parametrization by sex and smoking.} Let $z(t)\in\{\ns,\sm\}$
(non-smoker, smoker; IJM subscripts) indicate whether the individual accrues
smoking exposure at age $t$ (archival argument \code{smok\_status} of
\code{HT\_mal}, values \code{"N"}, \code{"Y"}):
\[
z(t) = \begin{cases} \sm & t_s \le t < t_q,\\ \ns & \text{otherwise.}\end{cases}
\]
$z$ depends on age and differs from the entry status $\zeta$: for never smokers
$z\equiv \ns$; for current smokers $z$ switches from $\ns$ to $\sm$ at $t_s$; for
former smokers $z$ is $\ns$, then $\sm$ on $[t_s,t_q)$, then $\ns$ again.

The non-smoker parameters $(\nu_\ns,\alpha_\ns,\gamma_\ns)$ depend on sex, and
$\mu_\ns=\nu_\ns$. For smokers (IJM Eqs.~2--3, Table~2),
\[
\kappa(q) = 1 + \alpha_1 q^{\alpha_2},\qquad
\alpha_\sm = \alpha_\ns\,\kappa(q),\qquad
\gamma_\sm = \gamma_\ns\,\kappa(q),\qquad
\mu_\sm = \mu_\ns,\qquad
\nu_\sm = \rho\,\nu_\ns,
\]
with sex-specific $(\alpha_1,\alpha_2,\rho)$. We write $\psi_\ns$ and
$\psi_\sm(q)$ for the two parameter sets (for a given sex). IJM has no symbol for
the initiation multiplier; we call it $\rho$. In the archival code, $\rho=1+p_4$,
$\alpha_1 = p_5$, $\alpha_2=p_6$, $(\alpha_\ns,\gamma_\ns,\nu_\ns) = (a_0,g_0,v_0)$.
Table~\ref{tab:tsce} compares the values. Where the sources disagree,
\code{fastMILC} uses the archival values.

\begin{table}[h]
\centering
\caption{TSCE parameters by source. IJM: Table~2 of \cite{ijm2017}, also used in
\code{data-raw/10\_baseline\_model\_parameters.R}. Archival:
\code{archival\_code/R/HT\_mal.R}. NZ: Table~A1 of \cite{nz2024}. The archival code
also lets $\alpha_\ns$ differ between the two parameter sets for males (7.7 for
$\psi_\ns$, 7.5 as the base of $\psi_\sm$).}\label{tab:tsce}
\begin{tabular}{@{}l l cccccc@{}}
\toprule
Sex & Source & $X$ & $\nu_\ns$ & $\alpha_\ns$ & $\gamma_\ns$ & $\rho$ & $(\alpha_1,\alpha_2)$ \\
\midrule
male   & IJM      & ``10e+7'' & $7.16\times10^{-8}$ & 7.7  & 0.090 & 1    & (0.6, 0.22) \\
male   & archival & $10^7$    & $7\times10^{-8}$    & 7.7 / 7.5 & 0.090 & 1.01 & (0.6, 0.22) \\
male   & NZ       & ``10e7'' & $7.16\times10^{-8}$ & 7.7  & 0.090 & 1    & (0.6, 0.22) \\
\midrule
female & IJM      & ``10e+7'' & $1.07\times10^{-7}$ & 15.82 & 0.071 & 0.98 & (0.5, 0.32) \\
female & archival & $10^7$    & $1\times10^{-7}$    & 15.5  & 0.071 & 1.02 & (0.5, 0.32) \\
female & NZ       & ``10e7'' & $1.07\times10^{-7}$ & 15.82 & 0.071 & 1.02 & (0.5, 0.32) \\
\bottomrule
\end{tabular}
\end{table}

\paragraph{Smoking history.} The hazard switches between parameter sets at
$t_s$ and $t_q$ and is evaluated on the attained-age clock in every segment
(the clock does not restart at $t_s$ or $t_q$):
\begin{equation}\label{eq:lambdaM}
\lambda_{\mathrm{mal}}(t) = h_{\psi_{z(t)}}(t) =
\begin{cases}
h_{\psi_\ns}(t) & 0 \le t < t_s,\\
h_{\psi_\sm(q)}(t) & t_s \le t < t_q,\\
h_{\psi_\ns}(t) & t \ge t_q,
\end{cases}
\qquad
\Lambda_{\mathrm{mal}}(t) = \sum_{j} \Big[H_{\psi_j}\big(\min(t,b_{j+1})\big) - H_{\psi_j}(b_j)\Big]_{+},
\end{equation}
where $0=b_0<b_1=t_s<b_2=t_q<b_3=\infty$, $\psi_j$ is the parameter set on
segment $j$, and the sum runs over segments with $b_j<t$. Never smokers have a
single segment; current smokers have two. After quitting, a former smoker has
the never-smoker hazard $h_{\psi_\ns}(t)$ at attained age; past smoking affects
$\Tmal$ only through the cumulative intensity accrued on $[t_s,t_q)$. This is
the archival hazard, and IJM describes the same piecewise construction. It
switches instantaneously at $t_s$ and $t_q$, whereas a TSCE model with
piecewise-constant parameters would carry the clonal population across the
switch and give a history-dependent hazard. We reproduce the archival form.

\subsection{Archival sampler}
\code{t\_mal} draws $u\sim\Unif(0,1)$ and solves
$\exp\{-\Lambda_{\mathrm{mal}}(t)\}=u$ with \code{uniroot} on $[0,150]$, using $u$
to select the segment in advance. If \code{uniroot} fails, it returns
$\Tmal=150$. This is an exact inverse-CDF sampler up to root-finding tolerance,
with a point mass at 150 for the $u$ whose root is beyond~150.

\paragraph{Property of the archival model.} $\Tmal$ is drawn from birth and is
\emph{not} conditioned on $\Tmal > \Tentry$. Individuals may enter with a
prevalent malignancy, and possibly with $\Tdiagn<\Tentry$.

\section{Tumor growth and disease progression}\label{sec:progression}

\subsection{Model}
Tumor volume $V(t)$ (mm$^3$) follows Gompertz growth (IJM Eq.~4) from the
volume of one malignant cell $V_0 = \pi d_0^3/6$, $d_0=0.01$\,mm, with scale
parameter $m$ and shape parameter $s = 31m$:
\[
\frac{V(t)}{V_0} = \exp\!\left\{\frac{s}{m}\left(1 - e^{-m t}\right)\right\},
\qquad V(t) = \frac{\pi}{6}\,d(t)^3 ,
\]
with $t$ in days since onset (IJM does not state the unit of $m$; the archival
code divides by 365). The time in years from onset to volume $v$ is the inverse
(\code{t\_vol}):
\begin{equation}\label{eq:tvol}
\tau(v) = -\frac{1}{365\,m}\log\!\left(1 - \frac{m}{s}\log\frac{v}{V_0}\right)
        = -\frac{1}{365\,m}\log\!\left(1 - \frac{1}{31}\log\frac{v}{V_0}\right),
\end{equation}
defined for $\log(v/V_0) < 31$. The maximum volume
$V_{\max}=\pi d_{\max}^3/6$, $d_{\max}=130$\,mm, gives $\log(V_{\max}/V_0)\approx 28.4$.

Three volume thresholds are random (IJM Table~1):
\[
\Vreg = 1000\,Z_{\mathrm{reg}},\quad \Vdist = 1000\,Z_{\mathrm{dist}},\quad
\Vdiagn = 1000\,Z_{\mathrm{diagn}},
\]
\[
Z_\bullet \sim \LN(\mu_\bullet, \mathrm{sd}_\bullet)\ \text{independently},
\]
truncated jointly to the set
\begin{equation}\label{eq:truncset}
\mathcal{A} = \{V_0 < \Vdiagn < V_{\max}\} \cap \{V_0 < \Vreg < \Vdist < V_{\max}\}.
\end{equation}
($Z$ is in ml; the factor 1000 converts to mm$^3$.) IJM constrains location and
scale to be equal, $\mu_\bullet = \mathrm{sd}_\bullet$ (Table~2; also
\cite{mdm2021}), so each log-normal has one free parameter. The subscripted
$\mu_\bullet$ are log-normal locations and are unrelated to the TSCE
conversion rate $\mu$; IJM uses the same letter for both. On $\mathcal{A}$,
$\Vdiagn$ is independent of $(\Vreg,\Vdist)$. Then
\[
\Treg = \Tmal + \tau(\Vreg),\qquad \Tdist = \Tmal + \tau(\Vdist),\qquad
\Tdiagn = \Tmal + \tau(\Vdiagn),
\]
\[
\mathrm{stage} = \begin{cases}\text{localized} & \Vdiagn < \Vreg,\\
\text{regional} & \Vreg \le \Vdiagn < \Vdist,\\
\text{distant} & \Vdiagn \ge \Vdist,\end{cases}
\qquad d_{\mathrm{diagn}} = \left(6\Vdiagn/\pi\right)^{1/3}\ \text{mm}.
\]
Since $\tau$ is increasing, the stage comparison on times and on volumes is the same.
The calibrated parameters are $(m,\mu_{\mathrm{reg}},\mu_{\mathrm{dist}},\mu_{\mathrm{diagn}})$,
with posterior medians $3.4\times10^{-4}$, 2.16, 5.62, 2.65 (IJM Table~2).

\subsection{Archival sampler and quirks}
\code{t\_prog} uses whole-vector rejection: it redraws all three volumes until
the draw falls in $\mathcal{A}$, which is an exact sampler of the truncated
distribution. One quirk:
\begin{enumerate}[nosep]
\item \textbf{Swapped parameter vectors.} The code draws $\Vreg$ with
  \code{cdist} and $\Vdist$ with \code{creg}. IJM defines
  $\Vreg\sim\LN(\mu_{\mathrm{reg}},\mathrm{sd}_{\mathrm{reg}})$ and
  $\Vdist\sim\LN(\mu_{\mathrm{dist}},\mathrm{sd}_{\mathrm{dist}})$. If the
  calibration in IJM and \cite{mdm2021} ran this code, the posterior labelled
  $\mu_{\mathrm{reg}}$ (median 2.16) was the location of $\Vdist$, and
  $\mu_{\mathrm{dist}}$ (5.62) the location of $\Vreg$, with the order
  constraint $\Vreg<\Vdist$ then truncating heavily. We have not been able to
  check which code the calibration used.
\end{enumerate}
The current \code{Population} class draws $\Vdist$ by a quantile shift
($u' = u_{\mathrm{reg}} + u(1-u_{\mathrm{reg}})$ applied to the $\Vdist$
quantile function). That shift gives the conditional law of $\Vdist$ given
$\Vdist > \Vreg$ only when $\Vreg$ and $\Vdist$ have the same distribution. It
also replaces truncation by setting out-of-bounds values to \code{NA}. Neither
matches \eqref{eq:truncset}.

\section{Death from other causes}\label{sec:oc}

\subsection{Tabulated input}
IJM estimates cumulative incidence functions (CIFs) for each cause from NHIS and
SEER data and denotes the estimates $\hat S$; we write $F_k$ (other causes) and
$G_k$ (lung cancer). For each stratum
$k_{\mathrm{other}} = (\ag(\Tentry), \sex, \zeta, \dg(q))$ (the intensity group is
used only for current smokers) the archival lists \code{never.other},
\code{former.other}, \code{current.other} give a right-continuous step function
$F_{k}(\tau)$, $\tau\in[0,\tau^{\max}_{k}]$, in years since entry, on knots
$0=\tau_{k,0}<\tau_{k,1}<\dots<\tau_{k,J_k}=\tau^{\max}_k$. The knots are
irregular: most spacings are 0.25 years, and spacings up to 5.5 years occur where
no deaths were observed. Horizons $\tau^{\max}_k$ are 16 years for never and
former smokers and 3.75--10 years for current smokers. All other-cause tables
have $F_k(0)=0$. The terminal values $F_k(\tau^{\max}_k)$ range from 0.006
(young never smokers) to 0.71 (old former smokers), and equal 1 in one stratum
(current smokers, ``>85'', female, ``21+'', horizon 3.75 years).

\subsection{Model (latent-time, as implemented)}
IJM samples age and cause of death ``by the inverse transform approach \ldots
based on CIF estimates''. The archival code does this by treating $\Tdo$ as an
independent latent time with marginal distribution
\begin{equation}\label{eq:TOC}
\Pr(\Tdo - \Tentry \le \tau) = F_{k_{\mathrm{other}}}(\tau), \quad \tau \le \tau^{\max}_k,
\end{equation}
followed by a tail rule for the residual mass $1-F_k(\tau^{\max}_k)$. Two
modeling facts:
\begin{itemize}[nosep]
\item The clock is time since entry, and the table is fixed by the age group
  at entry. A person entering at 41 uses the ``40--44'' table for the whole
  follow-up, also after age 45.
\item The tables are cumulative incidences in the presence of the competing
  cause. Using each as the marginal of an independent latent time is the
  archival modeling choice; a cumulative incidence and a net (latent)
  distribution differ. \code{fastMILC} keeps the archival interpretation.
\end{itemize}

\subsection{Archival sampler and quirks}
\code{tdeath\_other} draws $u_1,u_2\sim\Unif(0,1)$.
\begin{itemize}[nosep]
\item If $u_1 > F_k(\tau^{\max}_k)$ (tail): it sets $\ell = \Tentry + \tau^{\max}_k$,
  $r=110$, and returns $\Tdo = \Tentry + u_2(r-\ell)$. The intended rule was
  presumably $\ell + u_2(r-\ell)$, i.e.\ $\Unif(\Tentry+\tau^{\max}_k,\,110)$. As
  coded, $\Tdo\sim\Unif(\Tentry,\ 110-\tau^{\max}_k)$, which overlaps the
  tabulated window and has upper bound $110-\tau^{\max}_k$ (94 for 16-year
  tables).
\item Otherwise it picks the tabulated value of $F_k$ nearest to $u_1$ (by
  absolute difference, over \code{unique(est)}), takes the corresponding pair
  of adjacent knots from \code{unique(time)}, and returns $\Tentry$ plus a uniform
  draw between them. This approximates the inverse CDF with a
  piecewise-uniform density between knots. The nearest-value rule differs from
  the inverse-CDF rule $\min\{j: F_k(\tau_{k,j}) \ge u_1\}$.
\item Simulation of the lookup (40{,}000 draws per table) shows that the induced
  CDF lags the table by about one cell: for never smokers, ``80--84'', male, the
  archival CDF at 2.0 years is 0.076 against $F_k(2.0)=0.086$; the gap shrinks to
  under 0.002 after 8 years. For lung tables the atom is not reproduced (archival
  CDF at 0 is 0 against $G_k(0)=0.057$ for distant, ``60--64'', male).
\end{itemize}

\section{Death from lung cancer}\label{sec:lc}

\subsection{Tabulated input}
For each stratum $k_{\mathrm{lung}}=(\mathrm{stage}, \ag(\Tdiagn), \sex)$ the
list \code{ci.lung} gives a step function $G_k(\tau)$ in \emph{months} since
diagnosis, on irregular knots up to $\tau^{\max}_k \approx 222$--$251$ months.
Two features differ from the other-cause tables:
\begin{itemize}[nosep]
\item $G_k(0)>0$ for many strata (up to 0.137 for distant stage): an atom at
  diagnosis.
\item $G_k(\tau^{\max}_k)$ is large (up to 0.86), and in a few strata equals 1
  (e.g.\ the fallback stratum ``distant NA male''). One stratum has $G\equiv 0$.
\end{itemize}
Strata with age group ``NA'' exist because the archival \code{grep} lookup
falls through to them when $\ag(\cdot)$ returns \code{NA}.

\subsection{Model}
In the archival code, lung-cancer death can occur only after clinical diagnosis,
at any stage. $\Tdl$ is an independent latent time with
\begin{equation}\label{eq:TLC}
\Pr\!\big(12\,(\Tdl-\Tdiagn) \le \tau\big) = G_{k_{\mathrm{lung}}}(\tau),
\end{equation}
where the clock is time since diagnosis and the stratum is fixed by stage and
age group at diagnosis. $\Tdl$ is drawn for every individual, including those
with $\Tdiagn > \Tdo$; the minimum in Section~\ref{sec:assembly} removes those
draws. IJM (Section~2.1) states a different rule: a death is attributed to lung
cancer ``if and only if it occurs after transition to the distant state''. The
archival code does not impose this, and \code{fastMILC} follows the code
(Section~\ref{sec:discrepancies}).

\subsection{Archival sampler and quirks}
\code{tdeath\_lung} uses the same nearest-value lookup and within-interval
uniform draw as \code{tdeath\_other}, with times divided by 12. If
$u_1 > G_k(\tau^{\max}_k)$, it returns $\Tdl = \Tdiagn + \tau^{\max}_k/12$ with a
flag \code{event = 0}. The flag is not used in \code{nat\_hist}, so the
residual mass becomes a point mass of lung-cancer death at the horizon. The atom
$G_k(0)$ is handled through the nearest-value lookup between the knots at 0.

\section{Assembly}\label{sec:assembly}
\code{nat\_hist} combines the components:
\[
\Tdeath = \min(\Tdo, \Tdl),\qquad
\mathrm{cause} = \begin{cases}\text{other} & \Tdo < \Tdl\\ \text{lung} & \text{otherwise}\end{cases},
\]
\[
\Tpred = \Tentry + \Delta_{\mathrm{pred}},\quad
\mathrm{death} = \mathbf 1\{\Tdeath < \Tpred\},\quad
\Tfinal = \min(\Tdeath, \Tpred),
\]
\[
\mathrm{lung\_inc} = \mathbf 1\{\Tdiagn < \Tfinal\},\quad
\mathrm{excl} = \mathbf 1\{\Tdeath < \Tentry\}.
\]
$\mathrm{excl}=1$ flags individuals who died before entry. This happens through
$\Tdl$ when $\Tdiagn < \Tentry$ (Section~\ref{sec:onset}), and through the
archival other-cause tail when $\Tentry > 110-\tau^{\max}_k$ ($M-\tau^{\max}_k$
in \code{fastMILC}); the analyst excludes them.

\section{Discrepancies between IJM and the archival code}\label{sec:discrepancies}

\code{fastMILC} follows the archival code in every case below: the archival
number for items 1--4, and the archival structure for items 5--6.

\begin{enumerate}[nosep]
\item TSCE values (Table~\ref{tab:tsce}): $\nu_\ns$, $\alpha_\ns$ (female),
  the male smoker base $\alpha$ (7.5 vs 7.7), and $\rho$ differ. IJM gives
  $\rho = 0.98$ for females; the archival code and NZ give 1.02.
\item $X$: IJM writes ``10e+7'', which read literally is $10^8$; the code uses
  $10^7$. IJM probably means $1\mathrm{e}{+}7$.
\item $d_0$: IJM text says 0.01\,$\mu$m, IJM Table~1 and the code say 0.01\,mm.
\item IJM Table~1 prints the Gompertz law with $\exp(mt)$; IJM Eq.~4 and the code
  use $\exp(-mt)$.
\item Swapped \code{creg}/\code{cdist} in \code{t\_prog} (Section~\ref{sec:progression}).
\item Lung-cancer death: IJM requires transition to the distant state; the code
  requires diagnosis, at any stage, with stage-specific CIFs.
\end{enumerate}

\begin{thebibliography}{9}
\bibitem{ijm2017} S.~A. Chrysanthopoulou. MILC: A microsimulation model of the
natural history of lung cancer. \emph{International Journal of Microsimulation},
10(3):5--26, 2017.
\bibitem{milc2014} S.~A. Chrysanthopoulou. \emph{MILC: MIcrosimulation Lung Cancer
model}. R package version 1.0, 2014.
\bibitem{mdm2021} S.~A. Chrysanthopoulou, C.~M. Rutter, and C.~A. Gatsonis.
Bayesian versus empirical calibration of microsimulation models: a comparative
analysis. \emph{Medical Decision Making}, 41(6):714--726, 2021.
\bibitem{nz2024} E.~H. Nanayakkara, N.~Sibanda, and B.~Surya. A calibrated model
for lung cancer natural history: the case for New Zealand. \emph{Journal of
Simulation}, 2024. doi:10.1080/17477778.2024.2426629.
\end{thebibliography}

\end{document}
